\chapter{Метод}
\label{sec:method}

В этой главе вводятся обозначения и формализуются метрики, на основе
которых далее (глава~\ref{sec:experiments}) сравниваются классический
LLM-as-a-Judge, гибридные схемы (multi-prompt, consensus) и асессорская
разметка.

Всюду далее \emph{система} обозначает сравниваемую конфигурацию
(оцениваемую модель, версию промпта или схему LLM-судьи), а
\emph{конкурент} -- альтернативную конфигурацию, с которой проводится
сравнение (другая модель, предыдущая версия пайплайна или
классический LLM-as-a-Judge).
Названия конкретных моделей и продуктов в формулах не используются.

Общая схема предлагаемого пайплайна показана на
рисунке~\ref{fig:pipeline-overview}: на вход подаются запросы (из потока
или отобранной корзины), ответы сравниваемых систем $S$ и $C$ и
эталонная асессорская разметка; ядро -- множество судей $\set{J}_q$
(раздел~\ref{sec:method:hybrid}); на выходе -- оценки ответов, согласие с
асессорами и итоговый вердикт с обновлённой конфигурацией судей.

\begin{figure}[htbp]
    \centering
    \begin{tikzpicture}[
        font=\small,
        input/.style={draw, rounded corners=2pt, align=center,
            text width=27mm, minimum height=9mm, inner sep=3pt},
        core/.style={draw, very thick, rounded corners=3pt, align=center,
            text width=24mm, minimum height=22mm},
        output/.style={draw, rounded corners=2pt, align=center,
            text width=29mm, minimum height=8mm, inner sep=3pt},
        >={Stealth[scale=0.9]}, every edge/.style={draw, ->, thick}
    ]
        \node[core] (core) {Пайплайн\\LLM-as-a-Judge\\$\set{J}_q$};
        \node[input, left=16mm of core, yshift=13mm] (i1) {Запросы (поток / корзина)};
        \node[input, left=16mm of core] (i2) {Ответы систем $S$ и $C$};
        \node[input, left=16mm of core, yshift=-13mm] (i3) {Эталонная разметка (golden set)};
        \node[output, right=16mm of core, yshift=13mm] (o1) {Оценки ответов судьями};
        \node[output, right=16mm of core] (o2) {Согласие с асессорами};
        \node[output, right=16mm of core, yshift=-13mm] (o3) {Вердикт и новая конфигурация};
        \draw (i1) edge (core); \draw (i2) edge (core); \draw (i3) edge (core);
        \draw (core) edge (o1); \draw (core) edge (o2); \draw (core) edge (o3);
    \end{tikzpicture}
    \caption{Общая схема воспроизводимого пайплайна оценивания:
    входы (запросы, ответы сравниваемых систем, эталонная разметка),
    ядро (множество судей $\set{J}_q$) и выходы (оценки, согласие с
    асессорами, вердикт с обновлённой конфигурацией).}
    \label{fig:pipeline-overview}
\end{figure}

\section{Постановка задачи и обозначения}
\label{sec:method:notation}

Пусть $\set{Q}$ -- корзина запросов (диалогов), вынесенных на оценку,
$N = \abs{\set{Q}}$ -- её размер.
Каждый запрос $q \in \set{Q}$ оценивается множеством судей $\set{J}_q$ --
это могут быть асессоры (люди) либо LLM-судьи в рамках
multi-prompt/consensus схемы (раздел~\ref{sec:method:hybrid}).
Размер $\abs{\set{J}_q}$ называется \emph{перекрытием} запроса $q$.

Сравниваются две системы: $S$ (оцениваемая) и $C$ (конкурент).
Для запроса $q$ и судьи $a \in \set{J}_q$ через
$r_a(S, q), r_a(C, q) \in \{1, 2, 3\}$ обозначаются ранги ответов
$S$ и $C$, присвоенные судьёй $a$ (1 -- лучший ответ, 3 -- худший).

\paragraph{Чек-лист критериев и его применение к мультимодальным ответам.}
Ранг каждого ответа судья (асессор или LLM) выводит из чек-листа из
четырёх бинарных критериев, проверяемых \emph{строго по порядку}:
Relevance $\to$ Accuracy $\to$ Completeness $\to$ Readability.
Число пройденных чекбоксов $c \in \{0,1,2,3,4\}$ считается до первого
\texttt{false} (если критерий не пройден, последующие не оцениваются):
например, нерелевантный ответ получает $c = 0$ независимо от остального.
Поскольку ответ мультимодален (текст, изображения, таблицы, ссылки,
элементы интерфейса, раздел~\ref{sec:intro:motivation}), критерии
трактуются с учётом всех модальностей:
\begin{itemize}
    \item \textbf{Relevance} -- отвечает ли ответ \emph{в целом} на
        запрос. Пример нарушения: на запрос «покажи график цены» ответ
        даёт только текстовое описание без самого графика.
    \item \textbf{Accuracy} -- фактическая верность во всех модальностях:
        не только текст, но и числа в таблице, подписи к изображению,
        целевые страницы ссылок. Пример нарушения: таблица характеристик
        содержит неверное значение; ссылка ведёт не на заявленный
        источник.
    \item \textbf{Completeness} -- покрывает ли ответ запрос полностью, с
        нужным форматом и деталями. Пример нарушения: запрошена таблица
        сравнения, а дан сплошной текст; на изображении обрезана часть
        запрошенной информации.
    \item \textbf{Readability} -- структурированность и читаемость с
        учётом вёрстки: заголовки, списки, корректно отрисованные
        таблицы и превью. Пример нарушения: «сломанная» таблица или
        неотобразившееся изображение/превью.
\end{itemize}
Для многошаговых диалогов оценивается только \emph{последний} ответ, но
с учётом всей предыстории диалога; визуальная составляющая ответа при
автоматической оценке приводится к стабильному виду (скриншот ответа,
раздел~\ref{sec:experiments:pilot-multimodal}).

\paragraph{Визуально-дизайнерская ось оценки (C.R.A.P.).}
Визуально-дизайнерская ось (Contrast, Repetition, Alignment, Proximity) вводилась как экспериментальная метрика для оценки визуального качества ответов; результаты валидации приведены в разделе~\ref{sec:results:crap}.

\section{Поточечная метрика NPS}
\label{sec:method:nps}

Для абсолютной оценки качества системы $S$ на корзине $\set{Q}$ каждому
ответу присваивается вердикт $v(S, q) \in \{1, 2, 3\}$, где
$1$ -- хороший ответ, $2$ -- нейтральный, $3$ -- плохой.
Метрика NPS (Net Promoter Score) определяется как разность долей
хороших и плохих вердиктов:
\begin{equation}
    \label{eq:nps}
    \mathrm{NPS}(S) =
        \left(
            \frac{\abs{\{q \in \set{Q} : v(S,q) = 1\}}}{N}
            -
            \frac{\abs{\{q \in \set{Q} : v(S,q) = 3\}}}{N}
        \right)
        \times 100\%.
\end{equation}
\footnotesize Термин NPS используется по аналогии со структурой формулы (разность долей положительных и отрицательных оценок), но не связан с маркетинговой метрикой Net Promoter Score.\normalsize

Для сравнения двух систем используется разность их NPS:
\begin{equation}
    \label{eq:delta-nps}
    \Delta\mathrm{NPS}(S, C) = \mathrm{NPS}(S) - \mathrm{NPS}(C).
\end{equation}

Метрика \eqref{eq:nps} оценивает каждую систему независимо и не требует
совместного предъявления ответов $S$ и $C$ судье, поэтому используется
как офлайн-метрика мониторинга качества отдельной системы.

\section{Попарная метрика SbS-NPS}
\label{sec:method:sbs-nps}

Метрика \eqref{eq:delta-nps} чувствительна к калибровке шкалы
вердиктов $v$ у разных судей.
Для прямого сравнения двух систем на одном и том же запросе используется
попарная side-by-side метрика SbS-NPS, основанная на ранговом
голосовании.

\paragraph{Голос судьи.}
Для запроса $q$ и судьи $a \in \set{J}_q$ голос определяется сравнением
рангов \eqref{eq:vote}:
\begin{equation}
    \label{eq:vote}
    v_a(q) =
    \begin{cases}
        S, & r_a(S,q) < r_a(C,q), \\
        C, & r_a(C,q) < r_a(S,q), \\
        \mathrm{draw}, & r_a(S,q) = r_a(C,q).
    \end{cases}
\end{equation}

\paragraph{Подсчёт голосов и вердикт запроса.}
Обозначим через $W_S(q)$ и $W_C(q)$ число судей, проголосовавших за
$S$ и за $C$ соответственно:
\begin{equation}
    \label{eq:winvotes}
    W_S(q) = \abs{\{a \in \set{J}_q : v_a(q) = S\}}, \qquad
    W_C(q) = \abs{\{a \in \set{J}_q : v_a(q) = C\}}.
\end{equation}

Вердикт запроса определяется порогом $\theta$ (долей голосов, при
которой результат считается значимым; на практике
$\theta \approx 0.7 \cdot \abs{\set{J}_q}$):
\begin{equation}
    \label{eq:verdict}
    \mathrm{Verdict}(q) =
    \begin{cases}
        S, & W_S(q) \ge \theta \ \text{и} \ W_S(q) > W_C(q), \\
        C, & W_C(q) \ge \theta \ \text{и} \ W_C(q) > W_S(q), \\
        \mathrm{draw}, & \text{иначе}.
    \end{cases}
\end{equation}

Запросы с перекрытием $\abs{\set{J}_q} < \theta$ из корзины исключаются:
при таком перекрытии порог \eqref{eq:verdict} недостижим и вердикт
всегда был бы $\mathrm{draw}$, что искажает метрику.

\paragraph{Обоснование значения $\theta$.}
Коэффициент $0.7$ выбран как компромисс между чувствительностью и
устойчивостью вердикта и согласуется с двумя ориентирами из литературы
(глава~\ref{sec:related}).
Во-первых, он близок к практическому порогу надёжности PandaLM
($F1 \ge 85$--$90\%$, \cite{wang2023pandalm}): вердикт принимается лишь
при квалифицированном большинстве судей, что отсекает шумные расхождения.
Во-вторых, требование «значимого большинства, а не простого перевеса»
воспроизводит логику перевзвешивания по согласию из PRD~\cite{li2023prd}:
голоса судей-«одиночек» не определяют исход.
Конкретное значение в продуктовом цикле калибруется по golden set
(раздел~\ref{sec:method:data}): $\theta$ подбирается так, чтобы вердикты
схемы максимально согласовывались с асессорскими на эталонной выборке;
$0.7$ используется как разумное значение по умолчанию до калибровки.

\paragraph{Итоговая метрика.}
SbS-NPS -- это разность долей запросов, на которых победила система $S$
и на которых победил конкурент $C$:
\begin{equation}
    \label{eq:sbs-nps}
    \mathrm{SbS\text{-}NPS}(S, C) =
        \frac{
            \abs{\{q \in \set{Q} : \mathrm{Verdict}(q) = S\}}
            -
            \abs{\{q \in \set{Q} : \mathrm{Verdict}(q) = C\}}
        }{N}
        \times 100\%.
\end{equation}

Положительное значение \eqref{eq:sbs-nps} означает, что система $S$
в среднем предпочитается судьями чаще, чем конкурент $C$; нулевое --
что системы статистически неразличимы по качеству на корзине $\set{Q}$.

\section{Доверительный интервал и проверка гипотез}
\label{sec:method:bootstrap}

Поскольку \eqref{eq:sbs-nps} -- точечная оценка по конечной выборке
запросов, выводы о различиях формулируются как статистические гипотезы,
проверяемые bootstrap-методом.
Ниже вводятся bootstrap-оценка доверительного интервала и три решающих
правила: проверка смещения метрики относительно нуля (гипотеза H1),
проверка значимости сдвига между двумя замерами (гипотеза H2) и проверка
\emph{эквивалентности} двух замеров. Эти три правила и образуют процедуру
статистического доказательства эффективности схем
раздела~\ref{sec:method:hybrid}.

\subsection{Bootstrap-оценка доверительного интервала}
\label{sec:method:bootstrap:ci}

Из корзины $\set{Q}$ независимо извлекается $B$ выборок $\set{Q}_b$,
$b = 1, \dots, B$, размера $N$ с возвращением.
Для каждой выборки вычисляется $\mathrm{SbS\text{-}NPS}_b(S, C)$ по
формуле \eqref{eq:sbs-nps}.
Точечная оценка и границы $(1-\alpha)$-доверительного интервала
вычисляются как выборочное среднее и квантили эмпирического
распределения $\{\mathrm{SbS\text{-}NPS}_b\}_{b=1}^{B}$ уровней
$\alpha/2$ и $1 - \alpha/2$.
На практике используются $B = 10\,000$ и $\alpha = 0.05$ (95\%
доверительный интервал).

\paragraph{Алгоритм парного bootstrap.}
Алгоритм~\ref{alg:paired-bootstrap} формализует процедуру вычисления
доверительного интервала разности $\Delta = \theta(M_1) - \theta(M_2)$
парным bootstrap-методом.
Парность существенна: при вычислении обоих замеров на одной и той же
подвыборке $\set{Q}_b$ ковариация $\mathrm{Cov}(\theta_b(M_1),
\theta_b(M_2)) > 0$, что уменьшает дисперсию $\mathrm{Var}(\Delta_b)$
по сравнению с независимыми подвыборками и повышает мощность теста
(раздел~\ref{sec:method:bootstrap:h2}).

\begin{table}[htbp]
\centering
\caption{Алгоритм парного bootstrap для разности
$\Delta = \theta(M_1) - \theta(M_2)$}
\label{alg:paired-bootstrap}
\begin{tabular}{@{}p{0.85\textwidth}@{}}
\textbf{Вход:} корзина $\set{Q}$ размера $N$; функции замеров
$\theta(M_1)$, $\theta(M_2)$; число бутстрэп-итераций $B$;
уровень значимости $\alpha$ \\
\textbf{Выход:} $(1{-}\alpha)$-доверительный интервал для $\Delta$ \\[4pt]
\textbf{1.} Для $b = 1, \dots, B$ выполнить:
\begin{itemize}
    \item[1.1.] Сгенерировать подвыборку $\set{Q}_b$: для каждого
        $i \in \{1, \dots, N\}$ выбрать
        $q_i \sim \mathrm{Uniform}(\set{Q})$ с возвращением
    \item[1.2.] Вычислить $\theta_b(M_1) \leftarrow
        \mathrm{SbS\text{-}NPS}(M_1, \set{Q}_b)$ по формуле~\eqref{eq:sbs-nps}
    \item[1.3.] Вычислить $\theta_b(M_2) \leftarrow
        \mathrm{SbS\text{-}NPS}(M_2, \set{Q}_b)$ по формуле~\eqref{eq:sbs-nps}
    \item[1.4.] $\Delta_b \leftarrow \theta_b(M_1) - \theta_b(M_2)$
\end{itemize}
\textbf{2.} Отсортировать $\Delta_1, \dots, \Delta_B$ по возрастанию \\
\textbf{3.} $\mathrm{CI}_{1-\alpha}(\Delta) \leftarrow
    \bigl[\Delta_{(\lfloor \alpha B/2 \rfloor)},\;
          \Delta_{(\lceil (1 - \alpha/2) B \rceil)}\bigr]$
    \hfill\textit{(квантили эмпирического распределения)} \\
\textbf{4.} Вернуть $\mathrm{CI}_{1-\alpha}(\Delta)$
\end{tabular}
\end{table}

В алгоритме~\ref{alg:paired-bootstrap} подвыборки $\set{Q}_b$
извлекаются с возвращением (bootstrap), что обеспечивает
формальное допущение о независимости наблюдений внутри каждой
итерации~$b$ при сохранении парности замеров~$M_1$ и~$M_2$.
Порядковые статистики $\Delta_{(1)} \le \dots \le \Delta_{(B)}$
образуют эмпирическую функцию распределения разности~$\Delta$,
границы доверительного интервала -- квантили этого распределения.
Число итераций $B = 10\,000$ обеспечивает стабильность оценки
квантилей: стандартная ошибка оценки~97.5-го процентиля не
превышает $1/\sqrt{B} \approx 1\%$~\cite{efron1993bootstrap}.

\subsection{Гипотеза H1: смещение метрики относительно нуля}
\label{sec:method:bootstrap:h1}

Первый вопрос -- различимы ли системы $S$ и $C$ на корзине вообще:
\begin{equation}
    \label{eq:h1}
    H_0^{(1)}:\ \mathrm{SbS\text{-}NPS}(S, C) = 0
    \qquad \text{против} \qquad
    H_1^{(1)}:\ \mathrm{SbS\text{-}NPS}(S, C) \neq 0.
\end{equation}
$H_0^{(1)}$ отвергается на уровне $\alpha$, если доверительный интервал
из раздела~\ref{sec:method:bootstrap:ci} не содержит нуля; эквивалентно
вычисляется bootstrap-$p$-значение
$p = 2\min\!\big(\hat{P}(\mathrm{SbS\text{-}NPS}_b > 0),\,
\hat{P}(\mathrm{SbS\text{-}NPS}_b < 0)\big)$.
Это и есть проверка «смещения диффа NPS относительно нуля»: значимый
результат означает, что одна из систем устойчиво предпочитается судьями.

\subsection{Гипотеза H2: значимость сдвига между замерами}
\label{sec:method:bootstrap:h2}

Второй вопрос -- значимо ли \emph{отличаются} два замера метрики:
$M_1$ и $M_2$ (например, два прогона в разные моменты времени, либо две
схемы судьи, либо схема судьи и асессоры). Обозначим
$\theta(M) = \mathrm{SbS\text{-}NPS}$ по замеру $M$ и
\begin{equation}
    \label{eq:h2}
    \Delta = \theta(M_1) - \theta(M_2),
    \qquad
    H_0^{(2)}:\ \Delta = 0
    \quad \text{против} \quad
    H_1^{(2)}:\ \Delta \neq 0.
\end{equation}
Ключевой момент -- замеры \emph{парные}: оба считаются на одной и той же
корзине $\set{Q}$. Поэтому используется \emph{парный} bootstrap: на каждой
итерации $b$ извлекается одна выборка запросов $\set{Q}_b$, и по ней
вычисляются \emph{оба} замера $\theta_b(M_1)$, $\theta_b(M_2)$, после чего
$\Delta_b = \theta_b(M_1) - \theta_b(M_2)$.
Доверительный интервал строится по распределению $\{\Delta_b\}$.
Парность существенна: при общей корзине
$\mathrm{Var}(\Delta) = \mathrm{Var}(\theta_1) + \mathrm{Var}(\theta_2)
- 2\,\mathrm{Cov}(\theta_1, \theta_2)$, а положительная ковариация
(один и тот же набор запросов) уменьшает дисперсию разности и повышает
мощность по сравнению с трактовкой замеров как независимых.
$H_0^{(2)}$ отвергается, если $(1-\alpha)$-интервал $\{\Delta_b\}$ не
содержит нуля.

\subsection{Проверка эквивалентности: воспроизводит ли судья асессоров}
\label{sec:method:bootstrap:equiv}

Для валидации LLM-судьи цель \emph{обратная} H2: хороший судья должен
\emph{воспроизводить} замер асессоров, то есть давать \emph{нулевой}
сдвиг $\Delta = \theta(\text{судья}) - \theta(\text{асессоры})$.
Здесь «$H_0^{(2)}$ не отвергнута» (интервал $\Delta$ содержит нуль) само
по себе не является доказательством совпадения: отсутствие значимого
различия может объясняться недостатком мощности
(раздел~\ref{sec:experiments:power}).
Корректная процедура -- проверка эквивалентности (two one-sided tests,
TOST)~\cite{schuirmann1987tost} относительно практически незначимого
порога $\delta$ (margin эквивалентности).

\paragraph{Формализация TOST.}
Пусть $\delta > 0$ -- порог эквивалентности (margin of equivalence),
задаваемый как максимально допустимое отклонение SbS-NPS судьи от
асессорского замера (например, $\delta = 10$~п.п.,
раздел~\ref{sec:experiments:power}).
Процедура TOST эквивалентна проверке двух односторонних гипотез:
\begin{align}
    \label{eq:tost-lower}
    H_{0L}^{(\delta)}:\ &\Delta \le -\delta
        &\quad \text{против} \quad
        &H_{1L}^{(\delta)}:\ \Delta > -\delta, \\
    \label{eq:tost-upper}
    H_{0U}^{(\delta)}:\ &\Delta \ge +\delta
        &\quad \text{против} \quad
        &H_{1U}^{(\delta)}:\ \Delta < +\delta.
\end{align}
Эквивалентность принимается тогда и только тогда, когда \emph{обе}
нулевые гипотезы $H_{0L}^{(\delta)}$ и $H_{0U}^{(\delta)}$
отвергаются на уровне~$\alpha$:
\begin{equation}
    \label{eq:tost-decision}
    \text{Эквивалентность} \iff
        \bigl(\Delta_{\mathrm{lower}} > -\delta\bigr)
        \ \wedge\
        \bigl(\Delta_{\mathrm{upper}} < +\delta\bigr),
\end{equation}
где $\Delta_{\mathrm{lower}}$ и $\Delta_{\mathrm{upper}}$ --
нижняя и верхняя границы $(1{-}2\alpha)$-доверительного интервала
$\Delta$, вычисленного парным bootstrap-методом
(алгоритм~\ref{alg:paired-bootstrap}).
Уровень значимости TOST -- двусторонний~$\alpha$ (то есть каждый из двух
односторонних тестов проводится на уровне~$\alpha$), что
эквивалентно построению $(1{-}2\alpha)$-интервала и проверке его
вложенности в $(-\delta, +\delta)$.

Порог $\delta$ задаётся как practically significant difference и
определяется одним из двух способов: (а) на основе минимального
обнаружимого эффекта MDE (раздел~\ref{sec:experiments:power}) или
(б) на основе эмпирического разброса между асессорами на golden set.
В данной работе используется $\delta = 10$~п.п. по SbS-NPS --
величина, сопоставимая с типичным разбросом асессорских замеров
при перекрытии $|\set{J}_q| \in [7, 10]$
(раздел~\ref{sec:method:data}).
Этот порог опирается на меж-асессорскую дисперсию (inter-rater reliability): если асессоры сами расходятся во мнениях на $10$~п.п. при одинаковой разметке одних и тех же запросов, требовать от модели-судьи точности выше этого порога математически бессмысленно.

Только выполнение \eqref{eq:tost-decision} позволяет утверждать, что
схема судьи статистически \emph{эквивалентна} асессорам, а не просто
«неотличима из-за малой выборки».
Контринтуитивность TOST (нужно \emph{доказать} различие от нуля в
обе стороны, чтобы установить близость) объясняется тем, что
отсутствие отклонения $H_0$ стандартного теста~\eqref{eq:h2}
свидетельствует лишь о недостатке мощности, а не о совпадении --
именно это различие и фиксируется формализацией~\eqref{eq:tost-lower}
vs~\eqref{eq:tost-upper}.

\subsection{Решающие правила по этапам эксперимента}
\label{sec:method:bootstrap:rules}

Три правила применяются к этапам из таблицы~\ref{tab:experiments-design}
так:
\begin{itemize}
    \item \textbf{Этап 1} (классический судья vs асессоры): по H2
        \emph{ожидается значимый} сдвиг ($\Delta \neq 0$) -- это
        количественно доказывает смещённость одиночного судьи и задаёт
        точку отсчёта.
    \item \textbf{Этапы 2--3} (multi-prompt vs classic; consensus vs
        multi-prompt): эффективность метода доказывается тем, что он
        \emph{приближает} замер к асессорскому -- парным bootstrap
        проверяется значимость уменьшения $|\Delta_{\text{к асессорам}}|$
        (или роста согласия) при переходе к более сложной схеме.
    \item \textbf{Этап 4} (consensus vs асессоры): итоговое доказательство
        -- проверка \emph{эквивалентности} \eqref{eq:tost}: если
        $\mathrm{CI}(\Delta) \subset (-\delta, \delta)$, consensus-схема
        статистически воспроизводит асессорский вердикт.
\end{itemize}
Дополнительно проверяется согласованность \emph{решений} H1: судья
считается валидным, если по метрике он воспроизводит и знак, и значимость
$\mathrm{SbS\text{-}NPS}$, которые дали бы асессоры (то есть приходит к
тому же выводу о том, какая система лучше).

\section{Гибридный LLM-as-a-Judge: multi-prompt и consensus}
\label{sec:method:hybrid}

Метрики \eqref{eq:nps}--\eqref{eq:sbs-nps} определены через множество
голосов $\set{J}_q$ безотносительно к тому, кто эти голоса формирует --
асессоры или LLM-судьи.
В этом разделе формализуются три схемы формирования множества
$\set{J}_q$ при автоматической оценке: классическая (single-judge),
multi-prompt и consensus.
Все три схемы используют один и тот же механизм голосования
\eqref{eq:vote} и подсчёта \eqref{eq:winvotes}--\eqref{eq:verdict}, и
отличаются только тем, как формируются элементы $\set{J}_q$.

\subsection{Классическая схема (single-judge)}

В классической схеме каждый запрос $q$ оценивается один раз одной
парой (модель-судья, промпт):
\begin{equation}
    \label{eq:classic-judge}
    \set{J}_q^{\mathrm{classic}} = \{ a^* \}, \qquad
    \abs{\set{J}_q^{\mathrm{classic}}} = 1.
\end{equation}
При $\abs{\set{J}_q} = 1$ порог $\theta$ в \eqref{eq:verdict} достижим
только при $\theta \le 1$, то есть вердикт запроса полностью определяется
единственным голосом $v_{a^*}(q)$.
Это делает классическую схему чувствительной к стохастичности и
смещениям конкретной модели-судьи (раздел~\ref{sec:intro:motivation}):
ошибка единственного судьи напрямую переходит в ошибку вердикта.

\subsection{Multi-prompt: агрегация по формулировкам промпта}
\label{sec:method:hybrid:multiprompt}

Идея multi-prompt evaluation~\cite{mizrahi2024stateofwhatart} -- заменить
единственную формулировку промпта судьи набором $k$ семантически
эквивалентных формулировок $p_1, \dots, p_k$ и агрегировать их вердикты.

Пусть модель-судья $a^*$ с формулировкой промпта $p_i$ даёт голос
$v_{a^*, p_i}(q)$ по формуле \eqref{eq:vote}.
Множество судей для запроса $q$ в multi-prompt схеме -- это та же модель
$a^*$, применённая с $k$ разными промптами:
\begin{equation}
    \label{eq:multiprompt-judges}
    \set{J}_q^{\mathrm{mp}} = \{ (a^*, p_1), \dots, (a^*, p_k) \}, \qquad
    \abs{\set{J}_q^{\mathrm{mp}}} = k.
\end{equation}
Голос каждого элемента $(a^*, p_i) \in \set{J}_q^{\mathrm{mp}}$
определяется как $v_{a^*, p_i}(q)$ по формуле \eqref{eq:vote}, после чего
подсчёт голосов \eqref{eq:winvotes} и вердикт \eqref{eq:verdict}
вычисляются стандартным образом при $\abs{\set{J}_q} = k$.

Таким образом, multi-prompt схема не меняет механизм голосования, а
увеличивает перекрытие $\abs{\set{J}_q}$ с $1$ до $k$ за счёт
варьирования промпта при фиксированной модели-судье.
Это превращает чувствительность к промпту
(\cite{razavi2025benchmarking, arabzadeh2025humanai}) из
неконтролируемого источника ошибки в усредняемую компоненту голосования.

\paragraph{Генерация эквивалентных формулировок промпта.}
Набор $\{p_1, \dots, p_k\}$ должен состоять из семантически
эквивалентных, но текстуально различных формулировок одной и той же
задачи оценки. В работе используются три стратегии порождения:
\begin{enumerate}
    \item \emph{Перефразирование} -- одна и та же инструкция
        переписывается несколькими способами с сохранением смысла
        (компактная vs развёрнутая формулировка критериев).
    \item \emph{Варьирование глубины рассуждения} -- от «сразу выставь
        вердикт» до «сначала по каждому критерию опиши сильные/слабые
        стороны, затем вынеси вердикт» (перенос идеи EvalPlanner,
        раздел~\ref{sec:related:approaches}).
    \item \emph{Few-shot vs zero-shot} -- формулировки с $1$--$2$
        примерами разметки из тренировочной корзины и без примеров.
\end{enumerate}
Все формулировки проходят проверку на эквивалентность по тренировочной
корзине (раздел~\ref{sec:method:data}): формулировка принимается в
набор, если доля её расхождений с остальными по вердикту запроса не
превышает заранее заданного порога; это отсекает случайно неэквивалентные
переформулировки до их использования в голосовании.

\subsection{Consensus: агрегация по моделям-судьям}
\label{sec:method:hybrid:consensus}

Consensus-схема расширяет multi-prompt на несколько моделей-судей
$a_1, \dots, a_l$, каждая из которых применяется с собственным набором
из $m$ формулировок промпта:
\begin{equation}
    \label{eq:consensus-judges}
    \set{J}_q^{\mathrm{cons}} =
        \bigcup_{j=1}^{l} \{ (a_j, p_1^{(j)}), \dots, (a_j, p_m^{(j)}) \},
    \qquad
    \abs{\set{J}_q^{\mathrm{cons}}} = l \cdot m.
\end{equation}
Как и в \eqref{eq:multiprompt-judges}, каждый элемент
$(a_j, p_i^{(j)}) \in \set{J}_q^{\mathrm{cons}}$ даёт голос по формуле
\eqref{eq:vote}, а подсчёт голосов и вердикт вычисляются по
\eqref{eq:winvotes}--\eqref{eq:verdict} при $\abs{\set{J}_q} = l \cdot m$.

Consensus-схема -- упрощённый (без итеративного обмена ранжированиями)
вариант идеи PRD~\cite{li2023prd}: вместо последовательных раундов
дискуссии между судьями используется одноразовое голосование расширенного
множества $\set{J}_q^{\mathrm{cons}}$.
Такое упрощение снижает число обращений к моделям-судьям по сравнению с
итеративной дискуссией, но требует выбора $l$ и $m$ -- компромисс между
надёжностью (больше $\abs{\set{J}_q}$ -- устойчивее вердикт к
самопредпочтению отдельной модели~\cite{wataoka2024selfpreference}) и
стоимостью (раздел~\ref{sec:experiments:cost}).

\subsection{Сводная схема}

Три рассмотренные схемы образуют единую иерархию по числу голосов на
запрос:
\begin{equation}
    \label{eq:hierarchy}
    \abs{\set{J}_q^{\mathrm{classic}}} = 1
    \ \le\
    \abs{\set{J}_q^{\mathrm{mp}}} = k
    \ \le\
    \abs{\set{J}_q^{\mathrm{cons}}} = l \cdot m.
\end{equation}
Эта иерархия наглядно показана на рисунке~\ref{fig:scheme-hierarchy}:
от единственного голоса в классической схеме к $k$ голосам одной модели
с разными промптами и далее к $l \cdot m$ голосам нескольких моделей.

\begin{figure}[htbp]
    \centering
    \begin{tikzpicture}[
        font=\small,
        box/.style={draw, rounded corners=3pt,
            text width=33mm, minimum height=28mm, inner sep=4pt},
        vote/.style={circle, draw, fill=black!12, minimum size=3.2mm, inner sep=0pt},
        >={Stealth[scale=0.9]}
    ]
        \node[box] (a) {};
        \node[box, right=10mm of a] (b) {};
        \node[box, right=10mm of b] (c) {};

        \node[font=\bfseries] at ([yshift=-4.5mm]a.north) {Классическая};
        \node[font=\bfseries] at ([yshift=-4.5mm]b.north) {Multi-prompt};
        \node[font=\bfseries] at ([yshift=-4.5mm]c.north) {Consensus};
        \node at ([yshift=-10mm]a.north) {$\abs{\set{J}_q}=1$};
        \node at ([yshift=-10mm]b.north) {$\abs{\set{J}_q}=k$};
        \node at ([yshift=-10mm]c.north) {$\abs{\set{J}_q}=l\cdot m$};

        \node[vote] at ([yshift=6mm]a.south) {};
        \foreach \i in {0,1,2}
            \node[vote] at ([xshift=\i*5mm-5mm, yshift=6mm]b.south) {};
        \foreach \r in {0,1}
          \foreach \i in {0,1,2}
            \node[vote] at ([xshift=\i*5mm-5mm, yshift=4mm+\r*5mm]c.south) {};

        \draw[->, thick] (a) -- (b);
        \draw[->, thick] (b) -- (c);
        \node[below=1mm of a, font=\footnotesize, text width=33mm, align=center]
            {1 голос};
        \node[below=1mm of b, font=\footnotesize, text width=33mm, align=center]
            {$k$ промптов, 1 модель};
        \node[below=1mm of c, font=\footnotesize, text width=33mm, align=center]
            {$l$ моделей $\times$ $m$ промптов};
    \end{tikzpicture}
    \caption{Иерархия схем голосования \eqref{eq:hierarchy}: число судей
    (кружков-голосов) на запрос растёт от $1$ (классическая) к $k$
    (multi-prompt) и $l\cdot m$ (consensus); механизм голосования
    \eqref{eq:vote}--\eqref{eq:verdict} при этом не меняется.}
    \label{fig:scheme-hierarchy}
\end{figure}

В главе~\ref{sec:experiments} эти схемы сравниваются между собой и с
асессорской разметкой по метрике \eqref{eq:sbs-nps} с доверительным
интервалом (раздел~\ref{sec:method:bootstrap}), а в
главе~\ref{sec:results} -- по стоимости получения одного вердикта.

\section{Данные: корзины и эталонная разметка}
\label{sec:method:data}

Для применения метрик \eqref{eq:nps}--\eqref{eq:sbs-nps} и схем
голосования (раздел~\ref{sec:method:hybrid}) нужны два типа данных:
рабочая корзина запросов $\set{Q}$ и эталонная (golden set) разметка.

\textbf{Рабочая корзина} $\set{Q}$ -- выборка диалогов с мультимодальной
LLM-системой, на которых сравниваются система $S$ и конкурент $C$ (две
версии системы, либо система и классический LLM-судья, либо LLM-судья и
асессоры -- в зависимости от того, что сравнивается в конкретном
эксперименте, раздел~\ref{sec:experiments}).

На практике корзины формируются по одному из четырёх принципов отбора
запросов, различающихся ролью в пайплайне:
\begin{itemize}
    \item \textbf{поточная корзина} -- случайная выборка запросов из
        текущего трафика системы без балансировки по сценариям; отражает
        реальное распределение запросов и используется для регулярного
        мониторинга качества;
    \item \textbf{валидационная корзина} -- выборка, сбалансированная по
        типам сценариев (классам запросов), используется для сравнения
        версий judge-схем между собой на сопоставимых по составу
        подмножествах;
    \item \textbf{golden set} -- подмножество с эталонной асессорской
        разметкой (см. ниже), используется как точка отсчёта для
        калибровки и сравнения с асессорами;
    \item \textbf{тренировочная корзина} -- выборка, на которой подбираются
        формулировки промптов $p_1, \dots, p_M$ судей (раздел
        \ref{sec:method:hybrid:multiprompt}) и калибруется порог $\theta$;
        не пересекается с валидационной и golden set, чтобы избежать
        переобучения промптов на тестовых данных.
\end{itemize}
Запросы во всех четырёх корзинах -- диалоги пользователя с мультимодальной
LLM-системой: помимо текста реплик ответ может содержать изображения,
таблицы, ссылки и элементы интерфейса (раздел~\ref{sec:intro:motivation}),
что и обуславливает мультимодальность входа судей в
разделе~\ref{sec:experiments}.
Роли четырёх корзин и их связи показаны на
рисунке~\ref{fig:baskets}.

\begin{figure}[htbp]
    \centering
    \begin{tikzpicture}[
        font=\small,
        b/.style={draw, rounded corners=3pt, align=center,
            text width=34mm, minimum height=14mm, inner sep=3pt},
        c/.style={draw, very thick, circle, align=center,
            text width=15mm, inner sep=2pt, minimum size=20mm},
        >={Stealth[scale=0.85]}
    ]
        \node[c] (hub) {Виды корзин};
        \node[b, above left=10mm and 6mm of hub] (stream)
            {\textbf{Поточная}\\мониторинг качества продукта};
        \node[b, above right=10mm and 6mm of hub] (val)
            {\textbf{Валидационная}\\замер качества judge-схем};
        \node[b, below left=10mm and 6mm of hub] (golden)
            {\textbf{Golden set}\\эталон, валидация и калибровка $\theta$};
        \node[b, below right=10mm and 6mm of hub] (train)
            {\textbf{Тренировочная}\\подбор промптов $p_1,\dots,p_M$};
        \foreach \n in {stream, val, golden, train}
            \draw[-, gray] (hub) -- (\n);
    \end{tikzpicture}
    \caption{Четыре типа корзин и их роль в пайплайне: поточная
    (мониторинг), валидационная (сравнение схем), golden set (эталон и
    калибровка), тренировочная (подбор промптов). Корзины не
    пересекаются между собой во избежание переобучения и контаминации.}
    \label{fig:baskets}
\end{figure}

\textbf{Эталонная разметка (golden set)} -- подмножество $\set{Q}$ (или
отдельная выборка), для которого получены ранги $r_a(S,q)$, $r_a(C,q)$ от
асессоров-людей.
Golden set используется двояко:
\begin{itemize}
    \item как $\set{J}_q$ непосредственно -- для вычисления
        \eqref{eq:sbs-nps} в варианте «асессоры vs LLM-судья»
        (раздел~\ref{sec:experiments});
    \item как референс для калибровки порога $\theta$ в
        \eqref{eq:verdict} и для периодической валидации схем
        раздела~\ref{sec:method:hybrid} в продуктовом цикле.
\end{itemize}

\textbf{Перекрытие.}
Для запроса $q$ перекрытием $\abs{\set{J}_q}$ называется число судей
(асессоров или элементов схемы из раздела~\ref{sec:method:hybrid}),
давших голос по этому запросу.
Как отмечено в разделе~\ref{sec:method:sbs-nps}, запросы с
$\abs{\set{J}_q} < \theta$ исключаются из корзины при вычислении
\eqref{eq:sbs-nps}; на практике для асессорской разметки приемлемым
считается перекрытие $\abs{\set{J}_q} \in [7, 10]$.

\textbf{Порядки размеров корзин.}
Размеры корзин выбираются из баланса «статистическая значимость против
стоимости разметки» (раздел~\ref{sec:experiments:cost}): поточная
корзина -- порядка $500$ запросов на цикл (широкое покрытие трафика),
валидационная -- порядка $200$ (сбалансированно по сценариям),
golden set -- порядка $100$ запросов с дорогой экспертной разметкой,
тренировочная -- порядка $50$ для подбора промптов и калибровки $\theta$.
Поточная и валидационная корзины регулярно обновляются (порядка одного
обновления в неделю) для защиты от контаминации
(раздел~\ref{sec:related:reproducibility}); golden set обновляется реже
из-за стоимости экспертной разметки.
Эталонная выборка, используемая в экспериментах
главы~\ref{sec:experiments}, отбирается из golden set с фильтрацией по
наличию ответов обеих систем и валидного вердикта.

\section{Воспроизводимость пайплайна}
\label{sec:method:reproducibility}

Метрики \eqref{eq:sbs-nps} и доверительный интервал
(раздел~\ref{sec:method:bootstrap}) сопоставимы между прогонами только
если сам процесс получения вердиктов $\mathrm{Verdict}(q)$
воспроизводим.
Для этого пайплайн фиксирует четыре группы артефактов каждого прогона.

\begin{enumerate}
    \item \textbf{Версии.} Для каждого судьи $a \in \set{J}_q$
        фиксируется идентификатор версии модели и идентификатор/хэш
        формулировки промпта $p_i$.
        Это превращает пары $(a_j, p_i^{(j)})$ из
        \eqref{eq:consensus-judges} в воспроизводимые сущности: одна и
        та же пара при повторном запуске обращается к той же версии
        модели и тому же тексту промпта.
    \item \textbf{Параметры запуска.} Температура генерации, seed (где
        поддерживается API модели) и прочие параметры декодирования
        фиксируются для каждого судьи и сохраняются вместе с результатом.
    \item \textbf{Логи.} Для каждого запроса $q$ и судьи $a$ сохраняются
        входной промпт, полный ответ модели-судьи и извлечённый из него
        ранг $r_a(S,q), r_a(C,q)$ -- это позволяет постфактум
        пересчитать \eqref{eq:vote}--\eqref{eq:verdict} при изменении
        порога $\theta$ без повторного обращения к моделям.
    \item \textbf{Повторные прогоны (reruns).} Для оценки остаточной
        стохастичности конкретной конфигурации судьи $\set{J}_q$ прогон
        корзины $\set{Q}$ может выполняться многократно; разброс
        \eqref{eq:sbs-nps} между такими прогонами -- дополнительная (по
        отношению к bootstrap-CI, раздел~\ref{sec:method:bootstrap})
        оценка устойчивости результата.
\end{enumerate}

Совокупность этих артефактов позволяет воспроизвести
$\mathrm{Verdict}(q)$ и итоговую метрику \eqref{eq:sbs-nps} по
сохранённым логам без повторного запуска моделей-судей, а также
сравнивать схемы раздела~\ref{sec:method:hybrid} между собой при равных
условиях запуска.

\section{Архитектура программного решения}
\label{sec:method:architecture}

Реализация пайплайна организована в виде совокупности Python-модулей
(каталог \texttt{code/experiments/}, доступен в репозитории),
каждый из которых отвечает за один функциональный слой архитектуры,
показанной на рисунке~\ref{fig:architecture}.

\begin{figure}[htbp]
    \centering
    \begin{tikzpicture}[
        font=\small,
        layer/.style={draw, rounded corners=3pt, align=center,
            text width=40mm, minimum height=14mm, inner sep=3pt},
        >={Stealth[scale=0.9]}, every edge/.style={draw, ->, thick}
    ]
        \node[layer] (data) at (0, 0) {\textbf{Данные}\\
            cache\_v2.json, golden\_set\_full.xlsx};
        \node[layer, above=8mm of data] (judge) {\textbf{Судьи}\\
            faithful\_judge\_run.py, improve\_judge.py};
        \node[layer, above=8mm of judge] (agg) {\textbf{Агрегация}\\
            verdict(), scheme\_verdicts()};
        \node[layer, above=8mm of agg] (stat) {\textbf{Статистика}\\
            bootstrap(), paired\_gain(), TOST};
        \node[layer, above=8mm of stat] (viz) {\textbf{Визуализация}\\
            make\_result\_charts.py, analyze\_improve.py};
        \draw (data) edge (judge); \draw (judge) edge (agg);
        \draw (agg) edge (stat); \draw (stat) edge (viz);
    \end{tikzpicture}
    \caption{Архитектура программного решения: пять функциональных
    слоёв -- от данных до визуализации.}
    \label{fig:architecture}
\end{figure}

\textbf{Слой данных} (\texttt{cache\_v2.json}, \texttt{golden\_set\_full.xlsx}).
Файл \texttt{cache\_v2.json} содержит $100$ диалогов (50~train, 50~test),
извлечённых из golden set селекторным способом
(раздел~\ref{sec:experiments:pilot-multimodal}).
Каждая запись -- словарь с ключами \texttt{query} (текст запроса),
\texttt{ydx} (текст ответа системы $S$), \texttt{gpt} (текст ответа
конкурента~$C$), \texttt{verdict} (асессорский вердикт) и
\texttt{split} (train/test).
Запросы делятся на train и test стратифицированно: из каждого класса
вердиктов (draw, gpt\_win, ydx\_win) берётся фиксированная доля и
alternate-распределяется по сплитам (раздел~\ref{sec:experiments:setup}).

\textbf{Слой судей} (\texttt{faithful\_judge\_run.py},
\texttt{improve\_judge.py}).
Модуль \texttt{faithful\_judge\_run.py} реализует честный
ассессорский протокол: пер-ответные чекбоксы по порядку
(relevance$\to$accuracy$\to$completeness$\to$readability) со
«стопом» на первом false, независимая оценка ответов $A$ и $B$
(раздел~\ref{sec:method:notation}).
Вызов модели выполняется через HTTP API Ollama
(\texttt{localhost:11434/api/generate}) с параметрами
\texttt{temperature$=0.0$} и \texttt{num\_predict$=250$}.
Для каждого запроса позиция ответов $A/B$ рандомизируется
(seed определяется хэшем query\_id), что исключает position bias.
Промпты задаются в словаре \texttt{PROMPTS} с ключами
\texttt{c1\_compact}, \texttt{c2\_stoprule}, \texttt{c3\_reason}
-- три семантически эквивалентные формулировки
(раздел~\ref{sec:method:hybrid:multiprompt})
(полные тексты промптов приведены в Приложении~\ref{app:prompts}).
Результаты записываются в \texttt{votes\_faithful.jsonl} построчно:
каждая строка -- JSON-словарь с \texttt{qid}, \texttt{model},
\texttt{prompt}, \texttt{c\_s}/\texttt{c\_c} (число пройденных
чекбоксов), \texttt{s\_checks}/\texttt{c\_checks} (бинарные
массивы), \texttt{vote} (S/C/draw) и \texttt{sec} (время вызова).

\textbf{Слой агрегации} (\texttt{analyze\_staged.py}).
Функция \texttt{verdict\_from\_votes()} реализует пороговое
голосование \eqref{eq:verdict} при $\theta = 0.7$: вердикт
принимается при $\ge 70\%$ голосов за одну из систем.
Функция \texttt{scheme\_verdicts()} собирает множество голосов
$\set{J}_q$ для каждой схемы (classic, multiprompt, cons3, cons4)
по формулам~\eqref{eq:classic-judge}--\eqref{eq:consensus-judges}
и вычисляет вердикт запроса.

\textbf{Слой статистики} (\texttt{analyze\_staged.py}).
Функция \texttt{bootstrap()} реализует парный bootstrap
(алгоритм~\ref{alg:paired-bootstrap}): $B = 10\,000$ итераций,
квантильный доверительный интервал.
Функция \texttt{paired\_agreement\_gain()} вычисляет прирост
согласия с асессорами между двумя схемами (парный bootstrap).
Порог эквивалентности $\delta = 10$~п.п.
(раздел~\ref{sec:experiments:power}).

\textbf{Слой визуализации} (\texttt{make\_result\_charts.py},
\texttt{analyze\_improve.py}).
Скрипт \texttt{make\_result\_charts.py} генерирует четыре графика
для главы~\ref{sec:results}: SbS-NPS по схемам с bootstrap-CI
(рисунок~\ref{fig:results-sbs-stages}), матрицу согласованности
вердиктов (рисунок~\ref{fig:results-confusion}), bootstrap-
плотности (рисунок~\ref{fig:results-nps-density}) и кривую
Pareto «согласие vs стоимость» (рисунок~\ref{fig:results-pareto}).
Скрипт \texttt{analyze\_improve.py} генерирует графики
итеративных улучшений (раздел~\ref{sec:results:improve}).
